\documentclass[11pt]{article}
\usepackage[margin=1in]{geometry}
\usepackage{booktabs}
\usepackage{enumitem}
\usepackage{hyperref}
\usepackage{mdframed}
\usepackage{array}
\usepackage{longtable}
\usepackage{titlesec}
\usepackage{xcolor}

\definecolor{sectionbg}{RGB}{210,225,245}

\titleformat{\part}[block]{\normalsize\bfseries\centering}{}{0pt}{}
\titleformat{\section}{\normalsize\bfseries}{}{0pt}%
  {\colorbox{sectionbg}{\parbox{\dimexpr\linewidth-2\fboxsep\relax}{\strut#1\strut}}}
\titleformat{\subsection}{\normalsize\bfseries}{}{0pt}{}
\titleformat{\subsubsection}{\normalsize\bfseries}{}{0pt}{}
\titlespacing*{\part}{0pt}{1.5ex}{1ex}
\titlespacing*{\section}{0pt}{1.5ex}{0.5ex}
\titlespacing*{\subsection}{0pt}{1ex}{0.5ex}
\titlespacing*{\subsubsection}{0pt}{1ex}{0.3ex}

\newcommand{\concern}[1]{\textbf{Concern:} #1}
\newcommand{\ok}[1]{\textbf{OK:} #1}
\newcommand{\openq}[1]{\textbf{Question:} #1}

\begin{document}

%======================================================================
\part{Editor's Letter}
%======================================================================

%----------------------------------------------------------------------
\section{Editor Comment 1: Use Pedersen (2022) as Conceptual Motivation}
%----------------------------------------------------------------------

\subsection*{Request}
\begin{mdframed}
\textit{``The AE communicated to me that unless you are prepared to engage
more deeply with the model, it may be more effective to use Pedersen (2022)
as a conceptual motivation for empirical tests, rather than positioning the
paper as a direct test of the model. I completely agree. I think you should
emphasize more the empirical aspects of your study.''}
\end{mdframed}

\subsection*{What We Did}
Revised Abstract, Section~1, Section~2 (renamed ``Conceptual Framework and
Empirical Implications''), and Section~7 to downplay direct-test language
and use ``motivated by,'' ``inspired by,'' and ``drawing on'' instead.
Section~2 is concise: four short subsections (model setup, belief formation,
price formation, trading dynamics), each one paragraph before a hypothesis,
explicitly framed as motivation rather than testing.

\subsection*{Valerie's Comments}
\openq{Was Section~2 shortened relative to the previous submission? If so,
mention it in the response letter as further evidence of addressing this
comment.}

%----------------------------------------------------------------------
\section{Editor Comment 2: Differentiate Types of Social Media Interactions}
%----------------------------------------------------------------------

\subsection*{Request}
\begin{mdframed}
\textit{``Would it be possible to differentiate the type of social media
interactions? It makes a huge difference seeing an anonymous post versus
seeing Trump's tweets. You would need to uncover deeper insights and more
interesting patterns from the data to provide sufficient contribution to
the literature.''}
\end{mdframed}

\subsection*{What We Did}
Introduced a celebrity vs.\ other user split in Section~5.4 and Table~9,
defining celebrity users as those whose average PageRank exceeds the weekly
90th percentile. Celebrity tone has stronger effects on returns (rational
2~bps, fanatic 6~bps, na\"{i}ve 3~bps vs.\ 1~bp, 2~bps, 2~bps for other
users) and a particularly large retail flow effect for fanatic users (0.66\%
vs.\ 0.24\%).

\subsection*{Valerie's Comments}
\concern{Table~9 shows short-flow effects are entirely null for both celebrity
and other users (all Granger $p$-values $>57\%$), but the text does not
mention this. This is potentially inconsistent with the main results (Tables~5
and~7) where tone negatively predicts short flow in the high-influence
subsample. Needs at least one sentence of acknowledgment.}

\openq{One possible explanation: the celebrity split uses the whole sample,
not the high-influence subsample, so average influence is lower. A brief note
to this effect would preempt reviewer questions.}

%----------------------------------------------------------------------
\section{Editor Comment 3: Technological Progress and Competing Platforms}
%----------------------------------------------------------------------

\subsection*{Request}
\begin{mdframed}
\textit{``Can you relate the discussion to technological progress? Would
the emergence of competing platforms affect the results?''}
\end{mdframed}

\subsection*{What We Did}
Section~6.1 (Table~10): ChatGPT release (November~30, 2022) as cutoff ---
agent tones predict returns only pre-ChatGPT (RationalTone coefficient falls
from $0.0014^{***}$ to $0.0009$, insignificant). Section~6.2 (Table~11):
relative activity on StockTwits, X, and Seeking Alpha defines
platform-emergence periods --- belief transmission and return predictability
weaken during these periods.

A striking additional finding in Table~10 Panel~B: the short-flow relation
flips sign around ChatGPT. Pre-ChatGPT, short sellers back off from bullish
tone; post-ChatGPT, they trade against it. All six coefficients significant.

\subsection*{Valerie's Comments}
\concern{The text overstates attenuation in retail trading during
platform-emergence periods. Table~11 shows rational and na\"{i}ve tone still
significantly predict retail flow during emergence periods ($p=3.1\%$ and
$p=3.6\%$). Only fanatic tone is fully attenuated. The claim should be
qualified.}

\openq{The ChatGPT short-flow sign flip is one of the strongest results in
the paper. Consider elevating it in the abstract or introduction.}

\newpage
%======================================================================
\part{Reviewer 1}
%======================================================================

%----------------------------------------------------------------------
\section{R1 Comment 1: PVAR Insufficient to Document Belief Spillover}
%----------------------------------------------------------------------

\subsection*{Request}
\begin{mdframed}
\textit{``This strategy overlooks a key feature\ldots the impact of social
interactions on belief spillovers. Without properly considering the role of
social interactions, this testing approach struggles to distinguish
alternative explanations\ldots A rational investor in the current period
would plausibly switch to a na\"{i}ve or fanatic investor in the next period
due to changes in their posts. Given this dynamic reshuffling\ldots the PVAR
results may simply reflect the persistence or autocorrelation of individual
investors' beliefs rather than genuine belief spillovers.''}
\end{mdframed}

\subsection*{What We Did}
Replaced the rolling keyword-based classification with a GPT-5 LLM
classification using each user's complete posting history, making investor
type a stable, fixed label rather than a period-by-period assignment (2.96\%
rational, 5.86\% fanatic, 91.18\% na\"{i}ve). Section~4.1 is reframed to
clarify that the PVAR examines group-level belief spillovers, not individual
transitions.

\subsection*{Valerie's Comments}
\concern{The common-shocks alternative is never discussed or ruled out. Even
with stable type labels, the cross-group PVAR pattern could arise if all
investor types respond to the same news but rational investors react one week
earlier than na\"{i}ve investors --- with no social learning involved. The
high-influence concentration (Panel~B vs.\ Panel~C of Table~2) argues against
this, but this counter-argument is never made explicitly in the text.}

%----------------------------------------------------------------------
\section{R1 Comment 2: Investor Classification Ignores Fundamentals}
%----------------------------------------------------------------------

\subsection*{Request}
\begin{mdframed}
\textit{``One major concern is that this paper defines rational investors
solely based on view dynamics, without considering how these dynamics align
with the firm's fundamentals\ldots [Also:] the article defines na\"{i}ve
investors as those who demonstrate greater flexibility in adjusting their
views\ldots behavior that could indicate a higher level of sophistication.
This approach to classification diverges from the traditional definitions of
rational and na\"{i}ve investors in classical asset pricing theory.''}
\end{mdframed}

\subsection*{What We Did}
The LLM prompt includes the full Pedersen (2022) model description, and the
rational investor example (Encore Wire EBITDA analysis) shows the LLM rewards
fundamental-based reasoning. Earnings surprise validation (Appendix~B):
RationalTone $0.006^{**}$ ($t=2.12$), FanaticTone $-0.002$ ($t=-0.80$),
Na\"{i}veTone $-0.006^{**}$ ($t=-2.60$) --- only rational tone aligns with
fundamentals.

\subsection*{Valerie's Comments}
\concern{The conceptual critique --- that ``na\"{i}ve'' (flexible updaters)
could be viewed as more sophisticated than ``rational'' (stubborn) investors
in real markets --- is not addressed in the paper. The response letter
footnote acknowledges this tension but it does not appear in the manuscript.
Consider adding one sentence in Section~2 or~3.1.}

%----------------------------------------------------------------------
\section{R1 Comment 3: What Is Being Transmitted --- Information or Noise?}
%----------------------------------------------------------------------

\subsection*{Request}
\begin{mdframed}
\textit{``If this spillover primarily involves information, then rational
investors should, at the very least, exhibit predictive power comparable to
that of na\"{i}ve investors in forecasting future returns. However\ldots
rational investors are no better than na\"{i}ve investors at predicting
future return dynamics. This finding appears to contradict the `information
transmission' narrative.''}
\end{mdframed}

\subsection*{What We Did}
With the new LLM classification, rational tone now significantly predicts
returns with power comparable to na\"{i}ve tone: RationalTone $0.0012^{**}$
($t=2.53$), FanaticTone $0.0018^{***}$ ($t=2.80$), Na\"{i}veTone
$0.0012^{***}$ ($t=3.94$) (Table~3 whole sample). The short-run
predictability of na\"{i}ve tone is explained via an attention/sentiment
channel (Section~4.2).

\subsection*{Valerie's Comments}
Fully addressed. No further action needed.

%----------------------------------------------------------------------
\section{R1 Comment 4: Influence or Attention?}
%----------------------------------------------------------------------

\subsection*{Request}
\begin{mdframed}
\textit{``The Influence variable\ldots primarily captures the intensity of
investors' subsequent participation and does not immediately indicate
whether these follow-up commenters are genuinely `influenced' by the
original posters. Alternatively, this variable might instead serve as a
general proxy for investor attention.''}
\end{mdframed}

\subsection*{What We Did}
(1) Section~3.1 now states: ``We note that this measure may capture investor
attention as well as social influence.'' (2) PageRank-based centrality
introduced as an alternative measure (Section~5.2, Table~7), replicating
the main findings. (3) Correlation between centrality and attention is 0.25,
suggesting they capture distinct dimensions.

\subsection*{Valerie's Comments}
\concern{The correlation of 0.25 is reported as a bare number with no
specification (Pearson vs.\ Spearman) and no $p$-value. Consider adding:
``The Pearson (Spearman) correlation\ldots is 0.25 ($p < 0.01$).''}

%----------------------------------------------------------------------
\section{R1 Comment 5: Rational Investors Do Not Predict Trading Flows}
%----------------------------------------------------------------------

\subsection*{Request}
\begin{mdframed}
\textit{``Panel A of Table 4 shows that while the tone of na\"{i}ve
investors predicts trading flows, the tone of rational and fanatic investors
does not. This is puzzling, as the tone of rational and fanatic investors
influences na\"{i}ve investors.''}
\end{mdframed}

\subsection*{What We Did}
With the new classification, all three tones significantly predict retail
flows: RationalTone $0.0142^{***}$ ($t=6.80$), FanaticTone $0.0195^{***}$
($t=6.24$), Na\"{i}veTone $0.0087^{***}$ ($t=6.48$) (Table~4 Panel~A whole
sample).

\subsection*{Valerie's Comments}
Fully addressed. No further action needed.

%----------------------------------------------------------------------
\section{R1 Minor Comments}
%----------------------------------------------------------------------

\subsection*{Request}
Three minor comments: (1) typo ``insignificant'' vs.\ ``do not predict'';
(2) short sellers ``detecting influence'' overstated in abstract; (3) Reddit
users / retail investors overlap not acknowledged.

\subsection*{What We Did}
(1) Corrected in Section~4.4. (2) Toned down in abstract: ``Short sellers
\textit{may} account for agent influence\ldots This \textit{may} enhance
their ability\ldots'' (3) Acknowledged in Section~4.3.

\subsection*{Valerie's Comments}
Fully addressed. No further action needed.

\newpage
%======================================================================
\part{Reviewer 2}
%======================================================================

%----------------------------------------------------------------------
\section{R2 Comment C1: Mapping Between Empirical and Theoretical Agents}
%----------------------------------------------------------------------

\subsection*{Request}
\begin{mdframed}
\textit{``In the manuscript, both rational and fanatic agents are proxied
by Reddit users exhibiting a stable tone\ldots The empirical results suggest
that this proxy does not capture the distinct theoretical attributes\ldots
Na\"{i}ve tone is the strongest return predictor\ldots rational tone fails
to predict returns\ldots The authors could improve empirical alignment by
testing whether rational tone better predicts long-horizon returns or
fundamental outcomes (e.g., earnings surprises).''}
\end{mdframed}

\subsection*{What We Did}
(1) LLM reclassification resolves the return predictability issue --- all
three tones now comparable (see R1 Comment~3). (2) Earnings surprise test
(Appendix~B): only rational tone positively and significantly predicts
earnings surprises. (3) Long-horizon portfolio analysis (Section~4.5,
Figure~6): rational tone accumulates to $+4\%$ by week~52 (gradual price
discovery), na\"{i}ve tone reverses to $-1.6\%$ (transitory sentiment),
fanatic tone dissipates near zero ($+0.7\%$).

\subsection*{Valerie's Comments}
\concern{The fanatic tone dissipation ($+0.7\%$) is consistent with
Pedersen's bubble-unwind prediction --- fanatic bubbles eventually unwind
as rational agents reassert influence --- but this connection is never drawn
in the text. Add one sentence in Section~4.5 to complete the theoretical
narrative.}

%----------------------------------------------------------------------
\section{R2 Comment C2: Endogenous Influence Measure}
%----------------------------------------------------------------------

\subsection*{Request}
\begin{mdframed}
\textit{``The empirical proxy defines influence at the user $\times$ stock
$\times$ week level, based on the number of replies a user receives. This
operationalization conflates personal influence with transient engagement,
increasing the concerns of reverse causality\ldots The authors could isolate
exogenous components of influence: lagged cross-stock average influence,
pre-period network centrality, or decompose influence into fixed and residual
components.''}
\end{mdframed}

\subsection*{What We Did}
All three suggested robustness checks implemented: (1) lagged cross-stock
average commenters (Section~5.1, Table~6); (2) pre-period PageRank centrality
(Section~5.2, Table~7); (3) fixed-effect decomposition into persistent
user-level and transitory residual components (Section~5.3, Table~8). Main
findings are robust across all three. The negative tone--short-flow relation
is concentrated in the high user-fixed-effect subsample (Table~8 Panel~D),
suggesting it is driven by persistent rather than transitory influence.

\subsection*{Valerie's Comments}
\concern{Table~8 Panel~B shows return predictability remains statistically
significant even in the low-residual subsample (RationalTone $0.0010^{**}$,
FanaticTone $0.0014^{**}$, Na\"{i}veTone $0.0010^{***}$). The text claims
predictability is ``substantially stronger when either component is high''
but does not acknowledge the remaining significance in the low-residual group.
Add a qualifying sentence.}

%----------------------------------------------------------------------
\section{R2 Comment C3: Na\"{i}ve Investors and Retail Flow}
%----------------------------------------------------------------------

\subsection*{Request}
\begin{mdframed}
\textit{``Table 4 and Figure 4 reveal negligible and statistically
insignificant effects of either belief type on future retail flows across
all horizons\ldots Furthermore, the mechanism\ldots social media tone
$\to$ retail flow $\to$ price impact --- also lacks empirical support\ldots
The authors could test for Granger causality between tone and flow.''}
\end{mdframed}

\subsection*{What We Did}
Section~4.3 now explicitly acknowledges that retail flow should not be
interpreted as a direct proxy for na\"{i}ve investors' trading, addresses
the four concerns raised by R2 one by one, and reports bidirectional Granger
causality between tone and flow in the response letter: na\"{i}ve tone
$\to$ retail flow ($\chi^2=42.00$, $p=0.000$) and retail flow $\to$
na\"{i}ve tone ($\chi^2=47.72$, $p=0.000$).

\subsection*{Valerie's Comments}
\concern{The bidirectional Granger causality results appear in the response
letter only, not in the paper. Reviewers who look for this result in the
manuscript will not find it. These numbers should appear in a table or at
minimum in a footnote.}

%----------------------------------------------------------------------
\section{R2 Comment C4: Absence of Momentum-Reversal Pattern}
%----------------------------------------------------------------------

\subsection*{Request}
\begin{mdframed}
\textit{``Figures~3--5 display persistent price drifts following shocks to
fanatic/na\"{i}ve tone or retail flows, without evidence of mean reversion.
This contradicts Pedersen's model, which predicts reversal once rational
investors reassert influence.''}
\end{mdframed}

\subsection*{What We Did}
Section~4.5 and Figure~6: 52-week portfolio analysis sorting stocks on agent
tone signals shows the requested pattern --- na\"{i}ve tone reversal
($-1.6\%$ by week~52), rational tone continuation ($+4\%$), fanatic tone
dissipation ($+0.7\%$). This directly addresses the momentum-reversal
critique.

\subsection*{Valerie's Comments}
\concern{The reviewer asked for evidence that short sellers exert corrective
influence at longer horizons and for the full causal chain: tone $\to$ retail
flow $\to$ price drift $\to$ short flow $\to$ reversal. Figure~6 shows the
aggregate reversal but does not isolate the short-selling mechanism. Consider
either (a) adding a sentence acknowledging Figure~6 shows the aggregate
outcome of all corrective forces including but not limited to short selling,
or (b) adding short-flow-sorted long-horizon portfolios. Option (a) is the
lighter lift.}

\newpage
%======================================================================
\part{Issues to Fix}
%======================================================================

\section*{Must-Fix Before Resubmission}

\begin{enumerate}

  \item \textbf{Check footnote sequencing in the draft.} Footnotes jump
    from 5 to 9, then 11, then 14, skipping several numbers; and footnote
    14 appears twice. Check that no footnotes were accidentally dropped or
    duplicated during revision.

  \item \textbf{``P.\ Gal'' $\to$ ``P.\ Gao.''} The Da, Engelberg, Gao
    (2011) reference has the third author's name misspelled in the
    bibliography.

  \item \textbf{Atmaz (2024) missing from reference list.} Cited in
    Section~4.4 but absent from the reference list.

  \item \textbf{Engelberg, Reed, and Ringgenberg (2018) missing from
    reference list.} Cited in Section~4.4; the 2012 paper by the same
    authors is present but 2018 is a separate paper and needs its own entry.

  \item \textbf{Text--table inconsistency in Section~4.1.} The text says
    the RationalTone coefficient in the low-influence subsample (Panel~C) is
    $-0.004$; Table~2 Panel~C shows $-0.0004$. The text is off by one
    decimal place.

  \item \textbf{Table~9 note: ``other uses'' $\to$ ``other users.''}
    Typo in the table note.

  \item \textbf{``Commentors'' $\to$ ``commenters'' in Tables~6 and~8
    notes.} Appears in both table notes.

  \item \textbf{``Both'' $\to$ ``all three'' in Section~4.4.} The text
    states ``The corresponding p-values for the Granger causality tests are
    \textit{both} lower than 1\%'' when referring to rational, fanatic, and
    na\"{i}ve tone --- three variables.

  \item \textbf{Garbled math formula in Section~5.2 (before Equation~7).}
    The formula for the fixed-effect decomposition renders as a string of
    empty subscript/superscript delimiters. Verify it renders correctly in
    the final PDF.

  \item \textbf{Spurious capital ``F'' in Section~5.3.} The text reads
    ``For each component, \textbf{F}irms are divided\ldots'' --- should be
    lowercase ``firms.''

  \item \textbf{``Stefano, Corradin and Maddaloni (2020)'' missing from
    reference list.} Cited in footnote [10] but absent from the reference
    list.

  \item \textbf{Lyosca/Lyocsa: author name and year mismatch.} Footnote~[4]
    cites ``Lyocsa, Baumohl and Vyrost (2022)'' but the reference list has
    ``Lyosca \ldots 2021.'' Verify the correct spelling and year.

  \item \textbf{Several references may be unused in the main text.}
    Anderson et al.\ (2001), Blume and Stambaugh (1983), Box et al.\ (2021),
    Hilscher et al.\ (2015), Holtz-Eakin et al.\ (1988), Nickell (1981),
    Slezak (1994) appear in the reference list but are not cited in the main
    text or footnotes. Check appendices; if uncited, remove them.

\end{enumerate}

\section*{Substantive Issues for Coauthors}

\begin{enumerate}

  \item \textbf{Null short-flow result for celebrity users unreported.}
    Table~9 short-flow Granger $p$-values all $>57\%$ for both celebrity
    and other users. Needs at least one sentence of acknowledgment or
    explanation. (Editor Comment~2)

  \item \textbf{ChatGPT short-flow sign flip underemphasized.}
    All six pre/post-ChatGPT short-flow coefficients significant with a sign
    flip. Consider elevating in the abstract or introduction.
    (Editor Comment~3)

  \item \textbf{Text overstates retail-flow attenuation in Table~11.}
    Rational and na\"{i}ve tone still significantly predict retail flow during
    platform-emergence periods. Qualify the claim. (Editor Comment~3)

  \item \textbf{Common-shocks alternative not discussed.} Add one sentence
    using the high-influence concentration as the counter-argument.
    (R1 Comment~1)

  \item \textbf{Conceptual critique of type definitions not addressed in
    paper.} Add one sentence in Section~2 or~3.1. (R1 Comment~2)

  \item \textbf{Correlation between centrality and attention (0.25) not
    specified.} Add Pearson vs.\ Spearman and $p$-value. (R1 Comment~4)

  \item \textbf{Bidirectional tone--flow Granger causality in response
    letter only.} Add to paper as table or footnote. (R2 Comment~C3)

  \item \textbf{Table~8 Panel~B: low-residual return predictability.}
    Add a qualifying sentence acknowledging remaining significance.
    (R2 Comment~C2)

  \item \textbf{Fanatic tone dissipation not connected to Pedersen bubble
    mechanism.} Add one sentence in Section~4.5. (R2 Comment~C1)

  \item \textbf{Short sellers as corrective force at long horizons not
    demonstrated.} Acknowledge gap or add short-flow-sorted portfolios.
    (R2 Comment~C4)

\end{enumerate}

\end{document}
